\documentclass[aspectratio=169]{beamer}
\usetheme{Madrid}
\usepackage[spanish]{babel}
\usepackage{hyperref}
\usepackage{tikz}
\usetikzlibrary{arrows.meta,positioning}

\title{MarketPulse Beauty}
\subtitle{Retail Intelligence Dashboard para belleza y cuidado personal}
\author{Equipo MarketPulse Beauty}
\date{\today}

\begin{document}

\begin{frame}
  \titlepage
  \vspace{0.2cm}
  \begin{center}
    Fase 1: vender el valor \quad | \quad Fase 2: como se creo y como se ejecuto
  \end{center}
\end{frame}

\begin{frame}{Agenda}
  \begin{itemize}
    \item Fase 1: problema, oportunidad, propuesta de valor, demo conceptual
    \item Fase 2: datos, arquitectura, calidad, metodolog\'ia Scrum
    \item Cierre: resultados y siguientes pasos
  \end{itemize}
\end{frame}

\section{Fase 1: Vender el valor}

\begin{frame}{Por que esto importa a una pyme}
  \begin{itemize}
    \item Mucha data publica, poco tiempo para analizarla
    \item Reviews y tendencias cambian rapido y se detectan tarde
    \item Decisiones de producto se toman con intuicion, no con evidencia
  \end{itemize}
  \vspace{0.2cm}
  \textbf{Mensaje clave:} MarketPulse convierte ruido en decisiones claras.
\end{frame}

\begin{frame}{Propuesta de valor}
  \begin{itemize}
    \item Automatiza analisis de reviews y tendencias publicas
    \item Resume el estado del producto en indicadores simples
    \item Detecta alertas tempranas y oportunidades de mercado
    \item Reduce tiempo de analisis y acelera decisiones comerciales
  \end{itemize}
\end{frame}

\begin{frame}{Tres pilares del producto}
  \begin{itemize}
    \item Health Score del producto (rating + volumen de busquedas)
    \item Review Insights (temas recurrentes y sentimiento)
    \item Trend Analysis (interes de busqueda y picos de demanda)
  \end{itemize}
\end{frame}

\begin{frame}{Flujo de decision en el dashboard}
  \begin{center}
    \begin{tikzpicture}[
      node distance=1.6cm,
      box/.style={draw, rounded corners, align=center, minimum width=3.2cm, minimum height=0.9cm},
      arrow/.style={-Latex, thick}
    ]
      \node[box] (a) {Resumen ejecutivo};
      \node[box, right=of a] (b) {Cambios recientes};
      \node[box, right=of b] (c) {Motivo de cambios};
      \node[box, right=of c] (d) {Contexto de tendencias};
      \node[box, right=of d] (e) {Alertas y riesgos};
      \node[box, right=of e] (f) {Nivel de detalle};

      \draw[arrow] (a) -- (b);
      \draw[arrow] (b) -- (c);
      \draw[arrow] (c) -- (d);
      \draw[arrow] (d) -- (e);
      \draw[arrow] (e) -- (f);
    \end{tikzpicture}
  \end{center}
  \vspace{0.2cm}
  \textbf{Objetivo:} que el usuario entienda estado, causa y riesgo en minutos.
\end{frame}

\begin{frame}{KPIs y visuales que guian acciones}
  \begin{itemize}
    \item Rating promedio, Health Score, volumen de reviews
    \item Porcentaje de recomendacion y frecuencia de atributos
    \item Evolucion temporal de reviews y tendencias
    \item Visuales clave: ranking, barras horizontales, lineas temporales
  \end{itemize}
\end{frame}

\section{Fase 2: Como se construyo}

\begin{frame}{Como se construyo}
  \begin{itemize}
    \item Arquitectura orientada a datos con enfoque DataOps
    \item Pipeline de ingesta, limpieza, enriquecimiento y delivery
    \item Disenado para crecer: scraping, tendencias, alertas
  \end{itemize}
\end{frame}

\begin{frame}{Estrategia de datos}
  \begin{itemize}
    \item Fuente principal: Sephora (reviews ricas y con fecha)
    \item Arranque con dataset Kaggle (1M+ reviews, 8k+ productos)
    \item Plan de scraping incremental en sprints posteriores
    \item Limitaciones documentadas: sesgo de categorias, idioma, rating skew
  \end{itemize}
\end{frame}

\begin{frame}{Pipeline y arquitectura}
  \begin{itemize}
    \item Medallion: raw $\rightarrow$ bronze $\rightarrow$ silver $\rightarrow$ gold
    \item Orquestacion con Airflow (5 DAGs: ingestion, silver, gold, trends, data quality)
    \item Almacenamiento: MinIO (objetos) y PostgreSQL (gold layer)
    \item Dashboard en Streamlit, todo dockerizado
  \end{itemize}
\end{frame}

\begin{frame}{Inteligencia y calidad}
  \begin{itemize}
    \item Motor de insights por producto (rating, volumen, health score, trends)
    \item Reglas de tendencia con umbrales para alertas
    \item Monitoreo entre ejecuciones para detectar regresiones
    \item Validaciones de calidad para asegurar consistencia
  \end{itemize}
\end{frame}

\begin{frame}{Metodolog\'ia Scrum aplicada}
  \begin{itemize}
    \item 6 sprints de 2 semanas, con eventos formales
    \item Sprint 1: bases del repo, librer\'ias, fuentes de datos
    \item Sprint 6: producto anal\'itico, datasets derivados, Docker funcional
    \item Documentacion por sprint para trazabilidad y mejora continua
  \end{itemize}
\end{frame}

\section{Cierre}

\begin{frame}{Cierre: valor + ejecucion}
  \begin{itemize}
    \item MarketPulse traduce datos publicos en decisiones comerciales
    \item MVP con arquitectura escalable y metodolog\'ia solida
    \item Siguiente paso: demo con casos reales y validacion con pymes
  \end{itemize}
\end{frame}

\end{document}
